\documentclass[11pt]{article}

% ============================================================
% Packages
% ============================================================
\usepackage[utf8]{inputenc}
\usepackage[margin=1in]{geometry}
\usepackage{amsmath,amssymb,amsthm}
\usepackage{graphicx}
\usepackage{booktabs}
\usepackage{xcolor}
\usepackage{hyperref}
\usepackage{tikz}
\usepackage{enumitem}
\usepackage{pgfplots}

\pgfplotsset{compat=1.18}

% ============================================================
% Document Formatting
% ============================================================
\setlength{\parskip}{0.6em}
\setlength{\parindent}{0pt}

% ============================================================
% List Formatting
% ============================================================
\setlist[enumerate]{
    leftmargin=2em,
    itemsep=0.25em
}

% ============================================================
% Solution Environment
% ============================================================
\newenvironment{solution}
    {\par\noindent\textbf{Solution.}\par}
    {\par}

% ============================================================
% Link Styling
% ============================================================
\hypersetup{
    colorlinks=true,
    linkcolor=blue,
    urlcolor=cyan
}

% ============================================================
% Metadata
% ============================================================
\title{Ordinary Differential Equations}
\author{Chris Ritter}
\date{September 1, 2026}

% ============================================================
% Document
% ============================================================

\begin{document}

\maketitle
\tableofcontents
\newpage

% ============================================================
\section{First-Order Differential Equations}
% ============================================================

A first-order differential equation involves an unknown function and
its first derivative.

Some important types of first-order differential equations are:

\begin{itemize}
    \item Linear equations
    \item Separable equations
    \item Homogeneous equations
\end{itemize}

% ------------------------------------------------------------
\subsection{Linear Differential Equations}
% ------------------------------------------------------------

A first-order linear differential equation has the form

\[
\frac{dy}{dt} + p(t)y = g(t).
\]

Equivalently,

\[
\frac{dy}{dt} = -p(t)y + g(t).
\]

The standard method for solving a linear equation is the
\textbf{integrating factor method}.

The integrating factor is

\[
\mu(t)
=
e^{\int p(t)\,dt}.
\]

Multiplying the differential equation by \(\mu(t)\) gives

\[
\mu(t)y' + \mu(t)p(t)y
=
\mu(t)g(t).
\]

By construction,

\[
\frac{d}{dt}\left(\mu(t)y\right)
=
\mu(t)y' + \mu'(t)y
=
\mu(t)y' + \mu(t)p(t)y.
\]

Therefore,

\[
\frac{d}{dt}\left(\mu(t)y\right)
=
\mu(t)g(t).
\]

Integrating gives

\[
\mu(t)y
=
\int \mu(t)g(t)\,dt + C.
\]

Hence,

\[
\boxed{
y(t)
=
\frac{1}{\mu(t)}
\left(
\int \mu(t)g(t)\,dt + C
\right).
}
\]

% ------------------------------------------------------------
\subsection{Separable Differential Equations}
% ------------------------------------------------------------

A differential equation is \textbf{separable} if it can be written in
the form

\[
\frac{dy}{dx}
=
\frac{M(x)}{N(y)}.
\]

Rearranging gives

\[
N(y)\,dy = M(x)\,dx.
\]

Integrating both sides gives

\[
\int N(y)\,dy
=
\int M(x)\,dx.
\]

This generally produces an \textbf{implicit solution}.

Thus, the basic procedure is

\[
\boxed{
\frac{dy}{dx}
=
\frac{M(x)}{N(y)}
\quad\Longrightarrow\quad
N(y)\,dy=M(x)\,dx
\quad\Longrightarrow\quad
\int N(y)\,dy
=
\int M(x)\,dx.
}
\]

% ------------------------------------------------------------
\subsection{Homogeneous Differential Equations}
% ------------------------------------------------------------

A first-order differential equation is homogeneous if it can be
written in the form

\[
\frac{dy}{dx}
=
F\left(\frac{y}{x}\right).
\]

The standard substitution is

\[
v=\frac{y}{x}.
\]

Since

\[
y=xv,
\]

differentiating with respect to \(x\) gives

\[
\frac{dy}{dx}
=
v+x\frac{dv}{dx}.
\]

Substituting this into the original differential equation gives an
equation involving \(v\) and \(x\). This equation is often separable.

Thus,

\[
\boxed{
v=\frac{y}{x}
}
\]

transforms a first-order homogeneous equation into a separable
equation.

% ============================================================
\section{Modeling with First-Order Equations}
% ============================================================

Differential equations are frequently used to construct mathematical
models of real-world systems.

A general modeling process is:

\begin{enumerate}
    \item \textbf{Construct the model.}
    \item \textbf{Analyze the model.}
    \item \textbf{Compare the model with data and/or make predictions.}
\end{enumerate}

The first step is to identify the quantity being modeled.

For example, if \(A(t)\) represents the amount of some quantity at
time \(t\), then its rate of change can often be expressed as

\[
\boxed{
\frac{dA}{dt}
=
\text{sum of gain rates}
-
\text{sum of loss rates}.
}
\]

This principle is especially useful for population models and
mixing problems.

% ============================================================
\section{Population Model}
% ============================================================

\subsection{Problem}

An insect population in a region will increase at a rate proportional
to its current size. In the absence of any external factors, the
population triples in two weeks.

On any given day, there is a net migration of \(15\) insects into the
region, \(16\) insects are eaten by birds, and \(7\) insects die of
natural causes.

If initially there are \(100\) insects in the area, set up a
differential equation model to study whether this population will
survive. If not, estimate when the insects die out.

% ------------------------------------------------------------
\subsection{Step 1: Identify the Population}
% ------------------------------------------------------------

Let

\[
p(t)
\]

denote the insect population at time \(t\), measured in days.

The initial condition is

\[
\boxed{p(0)=100.}
\]

% ------------------------------------------------------------
\subsection{Step 2: Model the Natural Population Growth}
% ------------------------------------------------------------

The statement

\[
\text{``population increases at a rate proportional to its current size''}
\]

means

\[
\frac{dp}{dt}=rp,
\]

where

\[
r>0
\]

is a constant of proportionality.

This is the standard exponential growth model.

Separating variables gives

\[
\frac{1}{p}\,dp=r\,dt.
\]

Integrating,

\[
\int \frac{1}{p}\,dp
=
\int r\,dt,
\]

so

\[
\ln|p|=rt+C.
\]

Exponentiating,

\[
p=Ce^{rt}.
\]

Using \(p(0)=p_0\),

\[
p_0=Ce^0=C.
\]

Therefore,

\[
\boxed{
p(t)=p_0e^{rt}.
}
\]

% ------------------------------------------------------------
\subsection{Step 3: Determine the Growth Constant}
% ------------------------------------------------------------

In the absence of external factors, the population triples in
two weeks.

Since \(t\) is measured in days,

\[
p(14)=3p(0).
\]

Using the exponential model,

\[
p(14)=p_0e^{14r}.
\]

Therefore,

\[
p_0e^{14r}=3p_0.
\]

Canceling \(p_0\),

\[
e^{14r}=3.
\]

Taking the natural logarithm,

\[
14r=\ln(3).
\]

Thus,

\[
\boxed{
r=\frac{\ln(3)}{14}.
}
\]

% ------------------------------------------------------------
\subsection{Step 4: Account for External Factors}
% ------------------------------------------------------------

There are three external effects:

\begin{itemize}
    \item \(15\) insects enter the region per day.
    \item \(16\) insects are eaten by birds per day.
    \item \(7\) insects die naturally per day.
\end{itemize}

The total gain rate is

\[
15.
\]

The total loss rate is

\[
16+7=23.
\]

Therefore, the net external rate is

\[
15-23=-8.
\]

Thus,

\[
\boxed{
\frac{dp}{dt}=rp-8.
}
\]

Equivalently,

\[
\boxed{
\frac{dp}{dt}-rp=-8.
}
\]

This is a first-order linear differential equation.

% ------------------------------------------------------------
\subsection{Step 5: Solve the Differential Equation}
% ------------------------------------------------------------

The equation is

\[
p'-rp=-8.
\]

Comparing this with

\[
p'+a(t)p=g(t),
\]

we have

\[
a(t)=-r.
\]

Therefore, the integrating factor is

\[
\mu(t)
=
e^{\int -r\,dt}
=
e^{-rt}.
\]

Multiplying the differential equation by \(e^{-rt}\),

\[
e^{-rt}p'-re^{-rt}p
=
-8e^{-rt}.
\]

The left-hand side is

\[
\frac{d}{dt}
\left(
e^{-rt}p
\right).
\]

Therefore,

\[
\frac{d}{dt}
\left(
e^{-rt}p
\right)
=
-8e^{-rt}.
\]

Integrating,

\[
e^{-rt}p
=
\int -8e^{-rt}\,dt+C.
\]

Since

\[
\int -8e^{-rt}\,dt
=
\frac{8}{r}e^{-rt},
\]

we obtain

\[
e^{-rt}p
=
\frac{8}{r}e^{-rt}+C.
\]

Multiplying by \(e^{rt}\),

\[
p(t)
=
\frac{8}{r}
+
Ce^{rt}.
\]

Thus, the general solution is

\[
\boxed{
p(t)=\frac{8}{r}+Ce^{rt}.
}
\]

% ------------------------------------------------------------
\subsection{Step 6: Apply the Initial Condition}
% ------------------------------------------------------------

We know that

\[
p(0)=100.
\]

Therefore,

\[
100
=
\frac{8}{r}+C.
\]

Hence,

\[
C
=
100-\frac{8}{r}.
\]

Since

\[
r=\frac{\ln(3)}{14},
\]

we have

\[
\frac{8}{r}
=
\frac{8(14)}{\ln(3)}
=
\frac{112}{\ln(3)}
\approx 101.95.
\]

Therefore,

\[
C
\approx
100-101.95
=
-1.95.
\]

The population function is therefore

\[
\boxed{
p(t)
=
101.95
-
1.95e^{(\ln(3)/14)t}.
}
\]

% ------------------------------------------------------------
\subsection{Step 7: Analyze the Population}
% ------------------------------------------------------------

Because

\[
r=\frac{\ln(3)}{14}>0,
\]

we have

\[
e^{rt}\to\infty
\qquad\text{as }t\to\infty.
\]

Since the coefficient of the exponential term is negative,

\[
-1.95e^{rt}\to-\infty.
\]

Therefore,

\[
\boxed{
p(t)\to-\infty
\qquad\text{as }t\to\infty.
}
\]

Of course, a negative population is not physically meaningful.
Instead, the model tells us that the population reaches zero at
some finite time.

We therefore solve

\[
p(t^*)=0.
\]

Thus,

\[
101.95-1.95e^{(\ln(3)/14)t^*}=0.
\]

Rearranging,

\[
1.95e^{(\ln(3)/14)t^*}=101.95,
\]

so

\[
e^{(\ln(3)/14)t^*}
=
\frac{101.95}{1.95}.
\]

Taking the natural logarithm,

\[
\frac{\ln(3)}{14}t^*
=
\ln\left(\frac{101.95}{1.95}\right).
\]

Therefore,

\[
\boxed{
t^*
=
\frac{14}{\ln(3)}
\ln\left(\frac{101.95}{1.95}\right).
}
\]

Numerically,

\[
\boxed{
t^*\approx 40.2\text{ days}.
}
\]

Thus, according to this model, the population dies out after
approximately \(40\) days.

% ------------------------------------------------------------
\subsection{Equilibrium Population}
% ------------------------------------------------------------

An equilibrium population occurs when

\[
\frac{dp}{dt}=0.
\]

For

\[
\frac{dp}{dt}=rp-8,
\]

we set

\[
rp-8=0.
\]

Therefore,

\[
p^*=\frac{8}{r}.
\]

Using

\[
r=\frac{\ln(3)}{14},
\]

we obtain

\[
\boxed{
p^*
=
\frac{112}{\ln(3)}
\approx 101.95.
}
\]

This equilibrium is important because it determines the long-term
behavior of the model.

Since

\[
p(0)=100<101.95=p^*,
\]

the population begins below equilibrium.

Because \(r>0\), populations below the equilibrium decrease toward
zero and eventually reach zero in the physical interpretation of the
model.

% ============================================================
\section{Mixing Problems}
% ============================================================

Mixing problems model the amount of a substance in a container when
a solution enters and leaves the container.

The central idea is again

\[
\boxed{
\frac{dQ}{dt}
=
\text{rate in}
-
\text{rate out},
}
\]

where \(Q(t)\) is the amount of substance in the tank.

For a salt solution,

\[
\boxed{
\frac{dQ}{dt}
=
(\text{salt concentration in})
(\text{flow rate in})
-
(\text{salt concentration out})
(\text{flow rate out}).
}
\]

% ------------------------------------------------------------
\subsection{Salt Tank Example}
% ------------------------------------------------------------

At \(t=0\), a tank contains \(100\) gallons of water.

A solution containing

\[
\frac14\text{ lb of salt per gallon}
\]

enters the tank at a rate of \(r\) gallons per minute.

The tank is well stirred, and the solution leaves the tank at the
same rate \(r\) gallons per minute.

Let

\[
Q(t)
\]

denote the amount of salt in the tank at time \(t\), measured in
pounds.

Suppose

\[
Q(0)=Q_0.
\]

We want to:

\begin{enumerate}
    \item Find \(Q(t)\), the amount of salt in the tank at any time.
    \item Find the limiting amount of salt in the tank.
\end{enumerate}

% ------------------------------------------------------------
\subsection{Step 1: Determine the Rate of Salt Entering}
% ------------------------------------------------------------

The incoming solution has concentration

\[
\frac14\frac{\text{lb}}{\text{gal}}.
\]

The flow rate is

\[
r\frac{\text{gal}}{\text{min}}.
\]

Therefore,

\[
\text{rate in}
=
\left(
\frac14\frac{\text{lb}}{\text{gal}}
\right)
\left(
r\frac{\text{gal}}{\text{min}}
\right).
\]

The gallons cancel, leaving

\[
\boxed{
\text{rate in}=\frac{r}{4}\frac{\text{lb}}{\text{min}}.
}
\]

% ------------------------------------------------------------
\subsection{Step 2: Determine the Rate of Salt Leaving}
% ------------------------------------------------------------

Because the tank is well stirred, the concentration of salt leaving
the tank is the same as the concentration inside the tank.

At time \(t\), the tank contains \(Q(t)\) pounds of salt and
\(100\) gallons of solution.

Therefore, the concentration inside the tank is

\[
\boxed{
\frac{Q(t)}{100}
\frac{\text{lb}}{\text{gal}}.
}
\]

The solution leaves at a rate of

\[
r\frac{\text{gal}}{\text{min}}.
\]

Thus,

\[
\text{rate out}
=
\left(
\frac{Q(t)}{100}
\frac{\text{lb}}{\text{gal}}
\right)
\left(
r\frac{\text{gal}}{\text{min}}
\right).
\]

Therefore,

\[
\boxed{
\text{rate out}
=
\frac{rQ}{100}
\frac{\text{lb}}{\text{min}}.
}
\]

% ------------------------------------------------------------
\subsection{Step 3: Construct the Differential Equation}
% ------------------------------------------------------------

Using

\[
\frac{dQ}{dt}
=
\text{rate in}-\text{rate out},
\]

we obtain

\[
\frac{dQ}{dt}
=
\frac{r}{4}
-
\frac{rQ}{100}.
\]

Therefore,

\[
\boxed{
\frac{dQ}{dt}
=
\frac{r}{4}
-
\frac{r}{100}Q.
}
\]

Equivalently,

\[
\boxed{
\frac{dQ}{dt}
+
\frac{r}{100}Q
=
\frac{r}{4}.
}
\]

This is a first-order linear differential equation.

% ------------------------------------------------------------
\subsection{Step 4: Solve the Differential Equation}
% ------------------------------------------------------------

The equation has the form

\[
Q'+p(t)Q=g(t),
\]

where

\[
p(t)=\frac{r}{100}
\]

and

\[
g(t)=\frac{r}{4}.
\]

The integrating factor is

\[
\mu(t)
=
e^{\int (r/100)\,dt}
=
e^{rt/100}.
\]

Multiplying through by the integrating factor,

\[
e^{rt/100}Q'
+
\frac{r}{100}e^{rt/100}Q
=
\frac{r}{4}e^{rt/100}.
\]

The left-hand side becomes

\[
\frac{d}{dt}
\left(
e^{rt/100}Q
\right).
\]

Therefore,

\[
\frac{d}{dt}
\left(
e^{rt/100}Q
\right)
=
\frac{r}{4}e^{rt/100}.
\]

Integrating,

\[
e^{rt/100}Q
=
\int
\frac{r}{4}e^{rt/100}\,dt
+
C.
\]

Since

\[
\int
\frac{r}{4}e^{rt/100}\,dt
=
25e^{rt/100},
\]

we obtain

\[
e^{rt/100}Q
=
25e^{rt/100}+C.
\]

Thus,

\[
\boxed{
Q(t)=25+Ce^{-rt/100}.
}
\]

% ------------------------------------------------------------
\subsection{Step 5: Apply the Initial Condition}
% ------------------------------------------------------------

Since

\[
Q(0)=Q_0,
\]

we have

\[
Q_0=25+C.
\]

Therefore,

\[
C=Q_0-25.
\]

The solution is

\[
\boxed{
Q(t)
=
25+(Q_0-25)e^{-rt/100}.
}
\]

% ------------------------------------------------------------
\subsection{Step 6: Find the Limiting Amount}
% ------------------------------------------------------------

To find the limiting amount, consider

\[
\lim_{t\to\infty}Q(t).
\]

Since \(r>0\),

\[
e^{-rt/100}\to0
\qquad\text{as }t\to\infty.
\]

Therefore,

\[
\lim_{t\to\infty}Q(t)
=
25.
\]

Hence,

\[
\boxed{
\lim_{t\to\infty}Q(t)=25\text{ lb}.
}
\]

This makes sense physically because the incoming solution has
concentration

\[
\frac14\text{ lb/gal},
\]

and the tank contains \(100\) gallons. Therefore, at equilibrium,

\[
Q^*
=
\left(
\frac14\frac{\text{lb}}{\text{gal}}
\right)
(100\text{ gal})
=
25\text{ lb}.
\]

Thus, the equilibrium amount is exactly the amount corresponding
to the concentration of the incoming solution.

% ============================================================
\section{Key Modeling Principles}
% ============================================================

The examples above illustrate several important principles.

\subsection{General Balance Law}

For many models,

\[
\boxed{
\text{Rate of change}
=
\text{Rate in}
-
\text{Rate out}.
}
\]

More generally,

\[
\boxed{
\frac{dA}{dt}
=
\text{sum of gain rates}
-
\text{sum of loss rates}.
}
\]

\subsection{Population Models}

A population with growth proportional to its current size satisfies

\[
\boxed{
p'=rp.
}
\]

If there is also a constant net external contribution \(k\), then

\[
\boxed{
p'=rp+k.
}
\]

\subsection{Mixing Models}

For a well-stirred tank,

\[
\boxed{
Q'
=
(\text{incoming concentration})
(\text{incoming flow rate})
-
(\text{tank concentration})
(\text{outgoing flow rate}).
}
\]

The key assumption is that the tank is \textbf{well stirred}, so the
concentration of the material leaving the tank is equal to the
concentration throughout the tank.

\subsection{Equilibrium}

An equilibrium occurs when the quantity is not changing:

\[
\boxed{
\frac{dA}{dt}=0.
}
\]

For an equation

\[
A'=f(A),
\]

an equilibrium \(A^*\) satisfies

\[
\boxed{
f(A^*)=0.
}
\]

Equilibria are useful for understanding the long-term behavior of
solutions without necessarily solving the differential equation.

\end{document}
